
%\section{Introduction}
Sex trafficking forces individuals to perform sexual services against their will, often accompanied by an involuntary relocation of the exploited person.\footnote{For more information, see \href{https://www.oed.com/dictionary/sex-trafficking_n}{OED: Sex Trafficking}} Worldwide, there are estimated to be 27 million people exploited for labor, services, and sex.\footnote{For more information, see \href{https://www.state.gov/reports/2024-trafficking-in-persons-report/}{U.S. Department of State: Trafficking in Persons Report}} Sex traffickers target vulnerable and marginalized individuals, where women and girls account for 96\% of cases.\footnote{For more information, see \href{https://icat.un.org/sites/g/files/tmzbdl461/files/publications/icat-ib-04-v.1.pdf}{ICAT: The Gender Dimensions of Human Trafficking}} Traffickers exploit people who are economically disadvantaged and have little control over their economic outcomes (\textit{economic agency}). Factors that can increase sex trafficking risk include poverty, lack of social support, drug use, history of sexual abuse, being academically off-track, living in a shelter or group home, being an immigrant, or having a disability.\footnote{For more information, see \href{https://safesupportivelearning.ed.gov/human-trafficking-americas-schools/risk-factors}{U.S. Department of Education: Risk Factors}}  \\


Abortion access can affect women's \textit{economic agency}. As women have improved reproductive healthcare, they have more control over when they have children \citep{baileyandlindo}. This choice of being able to plan ahead allows for advances in employment. When women have children early in their careers, they reduce their working hours and firms invest less into their success \citep{BlauandKhan}. Then, young mothers are disadvantaged in the workforce and have less opportunity for promotions or job migration. In this sense, it is possible that as abortion laws become stricter, women lose economic agency, making them more at-risk for sex trafficking. 

Similarly, advancements in socioeconomic status often increases an individual’s exposure to information and civic engagement, increasing their awareness of social injustices.\footnote{For more information, see \href{https://www.urban.org/urban-wire/civic-engagement-higher-among-americans-who-are-financially-secure}{Urban Institute: Civic Engagement Is Higher among Americans Who Are Financially Secure}} As an individual becomes more aware of trafficking, they have greater knowledge to protect themselves therefore lowering their risk. Thus, reductions in economic agency through abortion access restrictions may decrease awareness and thus raise sex trafficking risk indirectly. 

To investigate this potential, I study the effect of the overturning of \textit{Roe v. Wade} on sex trafficking awareness. The period after \textit{Roe} is referred to as post-\textit{Dobbs}, as \textit{Dobbs v. Jackson Women’s Health Organization} (2022) is the case ruling that overturned \textit{Roe v. Wade} (1973).\footnote{For more information, see \href{https://reproductiverights.org/cases/scotus-mississippi-abortion-ban-dobbs-jackson-womens-health/}{ReproductiveRights.Org: Dobbs v. Jackson}}.\\
[A little more here on the general effect of the overturning of Roe on access]

In this paper, I ... more on what you actually did and found. 

Section 2 of this paper does this. Section 3 does that and so on. 

